\section{Filosofía de vida: largoplacismo, espíritu inicial y valentía}

\subsection{Los fundamentos matemáticos del largoplacismo}

Chen Tianqi (陈天奇) escribió en «Estos diez años de investigación en aprendizaje automático»: «Muchas de las dificultades y del desamparo no son más que parte de una fluctuación aleatoria; basta con invertir suficiente tiempo y paciencia para que el proceso aleatorio termine convergiendo a un estado estable acorde con lo invertido.»

Dice que esta frase «es literalmente un teorema»: el time average de un proceso aleatorio en efecto converge a una distribución estacionaria. Pero una cosa es conocer el teorema y otra muy distinta es perseverar en medio de los reveses.

\begin{importantbox}{La capacidad de aceptar el fracaso}
«El fracaso tampoco es tan terrible; después hay que aceptar el hecho de que uno puede fracasar.» Esta convicción se ha ido reforzando desde que en las Olimpiadas de Informática (OI) de la preparatoria solo obtuvo el segundo lugar, pasando por dos años sin resultados haciendo aprendizaje profundo durante la licenciatura, hasta los diversos giros de su emprendimiento. Chen Tianqi considera que: \textbf{mientras más fracasos se experimentan, más atrevimiento se tiene para actuar con soltura}, porque «fracasar tampoco es para tanto, de todos modos sigo viviendo bien». El fracaso temprano incluso es una especie de «protección»: si uno tiene demasiado éxito desde el principio, puede terminar generando dependencia de trayectoria.
\end{importantbox}

\subsection{La valentía proviene del yo del pasado}

Una idea profunda de Chen Tianqi es que \textbf{la valentía para actuar no aumenta automáticamente con la experiencia}. Por el contrario, mientras más cosas se ven, más fácil resulta volverse tímido y vacilante.

«Creo que muchas veces tu valentía tal vez sea el compromiso que tu yo del pasado hizo contigo mismo. Tienes que cumplir el compromiso que tu yo pasado hizo consigo mismo para poder seguir insistiendo en este style de trabajo.»

Él describe este proceso como un ciclo: el yo del pasado le da valentía al yo del presente, el yo del presente se anima continuamente a sí mismo durante el proceso de hacer las cosas, y luego le entrega esa valentía al yo del futuro.

\subsection{Mantener el espíritu inicial es lo más difícil}

Cuando le preguntaron «mirando hacia los próximos 5--10 años, ¿qué quieres hacer?», la respuesta de Chen Tianqi fue inesperada:

\textbf{«Creo que lo más difícil para mí ahora sigue siendo mantener el espíritu inicial.»}

Explica que, a medida que aumentan los recursos, aumentan las opciones y crece la responsabilidad, resulta más difícil mantener la actitud de aquellos años, la de «un becerro recién nacido no le teme al tigre». Cuando trabajaba en MXNet no le preocupaba en absoluto la competencia de TensorFlow, pero ahora, frente a un proyecto nuevo, no puede evitar preguntarse «si en realidad va a tener éxito o no». Mantener el espíritu inicial requiere un recordatorio deliberado hacia uno mismo.

El presentador insiste: los recursos que puedes movilizar son varios órdenes de magnitud mayores que hace 10 años, ¿por qué sientes que es más difícil mantener el espíritu inicial? Chen Tianqi responde: «Obtener recursos implica necesariamente responsabilidad. Tienes que garantizar que el equipo pueda tener éxito, tienes que producir el impact correspondiente para justify hacer las cosas relacionadas. Si te dejas arrastrar por estas presiones, pierdes la valentía para hacer cosas risky.»

Su estrategia para enfrentarlo es: seguir haciendo cosas más risky dentro del equipo de CMU, «estar dispuesto a renunciar y estar dispuesto a aceptar el fracaso, esta actitud en realidad es la misma que al principio». Los recursos y el equipo han crecido, y en efecto se pueden hacer más cosas que antes, pero el cambio de actitud puede hacer que uno «olvide esa actitud propia».

\begin{warningbox}{La trampa mental después del éxito}
Esta discusión revela una paradoja profunda: el éxito y la acumulación de experiencia pueden \textbf{erosionar la valentía necesaria para innovar}. Mientras más cosas se ven y más riesgos se conocen, más fácil resulta volverse conservador. Chen Tianqi considera que la manera de contrarrestar esta tendencia es volver al espíritu inicial: no tomar la decisión mediante un análisis racional de costos y beneficios, sino preguntarse a uno mismo: «¿es este el camino que idealmente quiero recorrer?». Si la respuesta es sí, entonces hacerlo sin importar el resultado.
\end{warningbox}

\subsection{Sobre el éxito y la felicidad}

En la parte final, el presentador hizo varias preguntas más personales.

\textbf{Sobre el éxito}: Chen Tianqi considera que el éxito tiene muchas definiciones: influencia, riqueza, avances tecnológicos, etcétera. Pero para él, lo más importante es «hacer cosas con la conciencia tranquila». «Poder influir en la gente en el ámbito técnico y poder contribuir en conjunto al mundo ya es algo bueno.»

\textbf{Sobre el mensaje para su yo futuro}: si le pidieran decirle una frase a su yo del pasado y otra a su yo del futuro, su elección no sería regresar a cambiar algo, sino todo lo contrario: \textbf{necesita que su yo del pasado le recuerde a su yo presente y futuro}: recordar el compromiso consigo mismo, sostener el ideal y seguir adelante.

\textbf{Sobre el estado ideal de vida}: el Chen Tianqi ideal de dentro de 10 o 20 años debería ser «un Chen Tianqi que ha tenido más fracasos, y que puede seguir manteniendo esa actitud inicial para desafiar lo que queremos desafiar».

\subsection{Resumen de la sección sobre la filosofía de vida}

La filosofía de vida de Chen Tianqi puede resumirse en tres palabras clave: \textbf{largoplacismo} (creer que el proceso aleatorio converge), \textbf{espíritu inicial} (hacer cosas con la conciencia tranquila) y \textbf{valentía} (aceptar el fracaso, sin dejarse arrastrar por los recursos y la responsabilidad). Estas no son consignas vacías, sino principios de acción verificados una y otra vez a lo largo de 20 años de investigación y emprendimiento. Vale la pena señalar que él no considera que estas cualidades se refuercen naturalmente con el aumento de la experiencia; por el contrario, considera que mantener el espíritu inicial y la valentía requiere un recordatorio constante y deliberado hacia uno mismo.
